Underwater acoustic (UWA) communication channels exhibit long delay spreads and 
severe fractional Doppler shifts that conventional orthogonal frequency division 
multiplexing (OFDM) cannot accommodate without prohibitive cyclic-prefix overhead. 
Affine Frequency Division Multiplexing (AFDM) folds these doubly-dispersive paths 
into a single sparse chirp-domain effective channel, enabling 
compressive-sensing-based channel estimation. To further preserve spectral 
efficiency in bandwidth-constrained UWA links, this work adopts a 
superimposed-pilot (SP) frame in which every chirp sub-carrier conveys both data 
and pilot energy. The resulting unknown-data interference does not decay with the 
noise variance, leaving a one-shot orthogonal matching pursuit (OMP) plus 
message-passing (MP) receiver estimation-limited at mid-to-high SNR.

This work proposes a complexity-preserving joint channel estimation and data detection
receiver for AFDM-SP on UWA-style sparse doubly-dispersive channels, in which OMP
is executed only once to lock the channel support and only a data-aided 
least-squares (LS) gain refit and sparse MP belief propagation iterate. An 
empirical support-recall probe motivates this design directly: the first OMP pass 
already recovers $\approx 95\%$ of the effective Dirichlet-expanded virtual-tap 
support, with the residual $\approx 5\%$ structurally determined by 
fractional-Doppler leakage geometry rather than by additive noise; re-running OMP 
across outer iterations is therefore a wasted constant-factor overhead. The 
proposed receiver thus avoids both this overhead and the cubic per-frame cost of a 
dense LMMSE estimator, while recovering two to three orders of magnitude in BER 
over the one-shot baseline at $\snr = 25$~dB. To recover the structural 
support-residual that OMP-once leaves behind, the receiver is further augmented 
with a 1\,041-parameter deep-unfolding support-refinement head that scores 
off-support candidates after the single OMP pass and appends a fixed number of 
top-ranked atoms. The learned augmentation stays in the same
$\mathcal{O}(\Ppath \Nsub^2)$ arithmetic class as the classical baseline and
incurs $\approx 8\%$ additional analytical complex-MAC overhead at
$\Nsub = 1024$, with measured GPU wall-clock parity at $\Nsub = 1024$
($402$~ms vs.\ $416$~ms for a classical-Torch reference) and a $+3.9\%$
wall-clock overhead at $\Nsub = 2048$.

A paired-frame Monte-Carlo evaluation on a UWA-style fractional-Doppler simulator 
confirms two consistent gains. The proposed classical joint receiver outperforms a 
serial OMP+MP baseline by two to three orders of magnitude in BER at $\snr = 
25$~dB. On top of that, the learned-augmented variant delivers a further $\approx 
2$~dB paired separation from the classical joint at $\snr \in \{20, 25, 30\}$~dB. 
A physical-support-oracle headroom analysis shows that the learned head closes 
$61\%/81\%/98\%$ of the achievable top-$\Ppath$ support-refinement gap at 
$\snr = 15/20/25$~dB, capturing the bulk of the BER reduction available from any 
support-refinement strategy operating on the same dictionary. Extension to 
multi-scale multi-lag UWA propagation and embedded-deployment latency, memory, 
and energy benchmarks are identified as natural follow-up work.